% from ega0-1-fr.tex
\oldpage[0\textsubscript{I}]{50}

identiques à $\mathcal{L}$~; on pose par convention
$\mathcal{L}^{\otimes0}=\mathcal{O}_X$ et pour $n\geq1$,
$\mathcal{L}^{\otimes(-n)}=(\mathcal{L}^{-1})^{\otimes n}$. Avec ces
notations, il existe alors un \emph{isomorphisme canonique fonctoriel}
\[
  \mathcal{L}^{\otimes m}\otimes_{\mathcal{O}_X}\mathcal{L}^{\otimes n}
  \xrightarrow{\sim}
  \mathcal{L}^{\otimes(n+m)}
  \tag{5.4.4.1}\label{0.5.4.4.1-fr}
\]
quels que soient les entiers rationnels $m,n$~: en effet, en vertu des
définitions, on se ramène aussitôt au cas où $m=-1$, $n=1$, et
l'isomorphisme en question a alors été défini en (5.4.3).
\end{env}

\begin{env}[5.4.5]
\EGAReaderTarget{canon-ref.0.5.4.5-fr}{0.5.4.5-fr}
Soit $f:Y\to X$ un morphisme d'espaces annelés. Si $\mathcal{L}$ est un
$\mathcal{O}_X$-Module localement libre (resp. inversible),
$f^*(\mathcal{L})$ est un $\mathcal{O}_Y$-Module localement libre (resp.
inversible)~: cela résulte aussitôt de ce que les images réciproques de deux
$\mathcal{O}_X$-Modules localement isomorphes sont localement isomorphes, de
ce que $f^*$ commute aux sommes directes finies et de ce que
$f^*(\mathcal{O}_X)=\mathcal{O}_Y$ (4.3.4). En outre, on sait qu'on a un
homomorphisme canonique fonctoriel
\[
  f^*(\check{\mathcal{L}})
  \longrightarrow
  \check{\bigl(f^*(\mathcal{L})\bigr)}
\]
(4.4.6), et lorsque $L$ est localement libre, cet homomorphisme est
\emph{bijectif}~: on est en effet encore ramené au cas où
$\mathcal{L}=\mathcal{O}_X$ qui est trivial. On en conclut que si
$\mathcal{L}$ est inversible, $f^*(\mathcal{L}^{\otimes n})$ s'identifie
canoniquement à $(f^*(\mathcal{L}))^{\otimes n}$ pour tout entier rationnel
$n$.
\end{env}

\begin{env}[5.4.6]
\label{0.5.4.6-fr}
Soit $\mathcal{L}$ un $\mathcal{O}_X$-Module inversible~; on désigne par
$\Gamma_*(X,\mathcal{L})$ ou simplement $\Gamma_*(\mathcal{L})$ le groupe
abélien somme directe
$\bigoplus_{n\in\mathbb{Z}}\Gamma(X,\mathcal{L}^{\otimes n})$~; on le munit
d'une structure d'\emph{anneau gradué}, en faisant correspondre au couple
$(s_n,s_m)$, où $s_n\in\Gamma(X,\mathcal{L}^{\otimes n})$,
$s_m\in\Gamma(X,\mathcal{L}^{\otimes m})$, la section de
$\mathcal{L}^{\otimes(n+m)}$ au-dessus de $X$ qui correspond canoniquement
(5.4.4.1) à la section $s_n\otimes s_m$ de
$\mathcal{L}^{\otimes n}\otimes_{\mathcal{O}_X}\mathcal{L}^{\otimes m}$~;
l'associativité de cette multiplication se vérifie de façon immédiate. Il est
clair que $\Gamma_*(X,\mathcal{L})$ est un foncteur covariant en
$\mathcal{L}$, à valeurs dans la catégorie des anneaux gradués.

Si maintenant $\mathcal{F}$ est un $\mathcal{O}_X$-Module quelconque, on pose
\[
  \Gamma_*(\mathcal{L},\mathcal{F})
  =\bigoplus_{n\in\mathbb{Z}}
   \Gamma\bigl(X,
     \mathcal{F}\otimes_{\mathcal{O}_X}\mathcal{L}^{\otimes n}\bigr).
\]
On munit ce groupe abélien d'une structure de \emph{module gradué} sur
l'anneau gradué $\Gamma_*(\mathcal{L})$ de la façon suivante~: au couple
$(s_n,u_m)$, où $s_n\in\Gamma(X,\mathcal{L}^{\otimes n})$ et
$u_m\in\Gamma(X,
\mathcal{F}\otimes_{\mathcal{O}_X}\mathcal{L}^{\otimes m})$, on fait
correspondre la section de
$\mathcal{F}\otimes_{\mathcal{O}_X}\mathcal{L}^{\otimes(m+n)}$ qui correspond
canoniquement (5.4.4.1) à $s_n\otimes u_m$~; la vérification des axiomes des
modules est immédiate. Pour $X$ et $\mathcal{L}$ fixés,
$\Gamma_*(\mathcal{L},\mathcal{F})$ est un foncteur covariant en
$\mathcal{F}$ à valeurs dans la catégorie des
$\Gamma_*(\mathcal{L})$-modules gradués~; pour $X$ et $\mathcal{F}$ fixés,
c'est un foncteur covariant en $\mathcal{L}$ à valeurs dans la catégorie des
groupes abéliens.

Si $f:Y\to X$ est un morphisme d'espaces annelés, l'homomorphisme canonique
(4.4.3.2)
\[
  \rho:\mathcal{L}^{\otimes n}
  \longrightarrow
  f_*\bigl(f^*(\mathcal{L}^{\otimes n})\bigr)
\]
définit un homomorphisme de groupes abéliens
\[
  \Gamma(X,\mathcal{L}^{\otimes n})
  \longrightarrow
  \Gamma\bigl(Y,f^*(\mathcal{L}^{\otimes n})\bigr),
\]
et comme
$f^*(\mathcal{L}^{\otimes n})=(f^*(\mathcal{L}))^{\otimes n}$, il résulte des
définitions des homomorphismes canoniques (4.4.3.2) et (5.4.4.1) que les
homomorphismes précédents définissent un \emph{homomorphisme fonctoriel
d'anneaux gradués}
$\Gamma_*(\mathcal{L})\to\Gamma_*(f^*(\mathcal{L}))$. Le même homomorphisme
canonique (4.4.3) définit de même un homomorphisme de groupes abéliens
\[
  \Gamma\bigl(X,
    \mathcal{F}\otimes_{\mathcal{O}_X}\mathcal{L}^{\otimes n}\bigr)
  \longrightarrow
  \Gamma\bigl(Y,
    f^*(\mathcal{F}\otimes_{\mathcal{O}_X}\mathcal{L}^{\otimes n})\bigr),
\]
et comme
\[
  f^*(\mathcal{F}\otimes_{\mathcal{O}_X}\mathcal{L}^{\otimes n})
  =f^*(\mathcal{F})\otimes_{\mathcal{O}_Y}
   (f^*(\mathcal{L}))^{\otimes n}
  \tag{4.3.3.1}
\]

